\section{Discussion}
\label{sec:discussion}

\subsection*{Clock-invariance and the fingerprint conjecture}
The Stage-1 results in Table~\ref{tab:stage_one_rqa} and
Figure~\ref{fig:rqa_violins} bear on the central conjecture of
Section~\ref{sec:problem}: that recurrence-plot geometry is a
fingerprint of the underlying microstructure. With two cells reported
in this run, the conjecture cannot be tested as a falsifiable
cross-venue claim (the $H_1$ Spearman test of
Equation~\ref{eq:h1} requires $M \geq 5$ instruments under the
$B = 10{,}000$ label-permutation null to admit a one-sided
$p \leq 0.05$). What we can report is a descriptive pattern: the
spread of $\overline{DET}$ \emph{across cells} within a fixed clock
is comparable to the spread \emph{across clocks} within a fixed cell.
For physical clock, binance and kraken sit at $\overline{DET} \in
\{0.626, 0.494\}$; for info\_time, $\{0.801, 0.615\}$. The
within-clock spread is on the order of $\Delta \overline{DET} \sim
0.1$--$0.2$, and the within-cell spread across clocks is on the same
order. This is the first sign that the clock choice is not a free
parameter --- it materially changes the recurrence-geometry summary
of the same underlying LOB --- and it argues for reporting every
downstream finding under all three clocks, as our pipeline already
does.

The info\_time clock concentrates mass at high determinism and
laminarity in both cells. Mechanically, this reflects the
volume-thresholded bucketer's design: each info\_time bucket spans a
calibrated $30$--$120$ second window in which a fixed cumulative
volume is exchanged, so the within-bucket trajectory tends to be
locally stationary by construction. The mean-recurrence-plot panels
in Figure~\ref{fig:mean_rp} make this visible: the info\_time panels
exhibit broad off-diagonal block structure characteristic of long
laminar segments, whereas the event panels concentrate mass near the
main diagonal with shorter diagonal lines. This is consistent with
article-1's qualitative reading of the four RP regimes
(deterministic, stochastic, stationary, oscillatory) and supports
treating the three clocks as distinct distributions rather than
redundant views of the same underlying process.

\subsection*{Stage-2 training behaviour and remaining unmeasured
axes}
Four of the six target (cell, clock) trainings completed at least the
first checkpoint (step $1{,}000$), as shown in
Figure~\ref{fig:em_gap}. EM-gap final-100-step means span
$[17.3, 51.2]$ across the four trained pairs; the spread is
$\sim 3\times$ larger than the per-pair final-100-step standard
deviation, suggesting that within-trajectory noise at this short
horizon is much smaller than cross-(cell, clock) variation in
asymptotic Wasserstein distance. Two trainings (binance/event and
binance/info\_time) did not produce stage-2 artefacts in this run;
the binance/info\_time pool of $40$ buckets sits at the
batch-size-64 viability boundary and is most likely the reason it
was not scheduled. The natural remediation is the production-grid
re-run discussed in the
introduction, not a change to the architecture.

The locked Stability Score depends on mode-coverage and critic-loss
curvature, which are populated only at $\geq 5{,}000$
generator-step checkpoints; in this run training stopped at the
first checkpoint and both axes are empty. The 2-axis partial Score
we report in Section~\ref{sec:results} is an ordering across the four
trained pairs, not a value comparable to article-1's locked
4-axis aggregate. The two findings worth carrying forward are: (i)
the kraken cell is the more stable training target on this short
horizon, with all three kraken clocks ranking above binance/physical
on the 2-axis partial score; (ii) the GP-magnitude across all four
trained pairs sits well above the Lipschitz target of $1$ (the four
mean GP values are $\{1.22, 1.36, 1.58, 1.65\}$), indicating that the
locked $\lambda_{\mathrm{GP}} = 10$ is too low to enforce the
constraint at the present generator capacity at step $1{,}000$. This
matches article-1's observation that a downward sweep of
$\lambda_{\mathrm{GP}}$ rather than corpus expansion is the natural
follow-on, and we adopt that framing for the production grid.

\subsection*{Limitations}
Three limitations are explicit. First, only two cells are reported,
which precludes the cross-venue $H_1$ Spearman test in any
statistically meaningful form. The natural minimum population is
$M = 5$ cells. Second, the seed budget for the reported stage-2
trainings is $J = 1$ (only \texttt{seed\_idx} $= 0$ exists on disk
for every trained pair); per-cell variance, the empirical primitive
of the Stability Score, is therefore not estimable from this run.
The natural minimum is $J = 5$. Third, the training horizon of
$1{,}000$ generator steps is too short to populate the
mode-coverage and curvature axes of the locked Stability Score;
extending the horizon to $\geq 5{,}000$ generator steps with
checkpointing at that cadence is a precondition for reporting the
full locked aggregate. None of these limitations changes the
architecture under study; all three are scaling and seed-population
questions that we treat as the natural follow-on path to the
production grid.
